\chapter{Concluzii}
\label{cap:concluzii}

\section{Concluzii privind realizarea temei}

Lucrarea a realizat o platformă coerentă, centrată pe integritatea datelor, pentru antrenamentele de
forță: planurile sunt transformate în istoric imuabil prin instantanee, analizele și feedback-ul sunt
derivate onest din date reale, iar un marketplace permite partajarea și copierea independentă a
programelor. Ingineria a pus accent pe corectitudine (teste pe baze de date reale), integritatea
datelor (instantanee pe două niveluri, ștergere logică, constrângeri la nivel de bază de date),
securitate (autentificare JWT cu rotație și acces sigur față de IDOR, RBAC plus ReBAC pentru coaching)
și explicabilitate (un analizor determinist, însoțit de avertismente), în locul unor funcționalități
superficiale.

Componenta de căutare a fost tratată ca o problemă de inginerie a sistemelor: un model de citire
derivat, eventual consistent și reconstruibil, cu securitate care nu poate fi ocolită din parametrii
cererii. Cea mai importantă contribuție din punct de vedere științific este \emph{evaluarea empirică}
OpenSearch vs.\ PostgreSQL pe un corpus scalat până la două milioane de documente, care a identificat
un punct de echilibru în jurul a 10.000 de documente și a arătat că diferența decisivă apare la
căutarea tolerantă la greșeli, unde abordarea relațională se degradează la mai multe secunde, în timp
ce motorul dedicat rămâne practic constant. Astfel, o decizie de arhitectură luată adesea „din
reflex” a fost transformată într-una informată și măsurabilă.

\section{Aprecieri personale privind relevanța rezultatelor}

Consider că valoarea lucrării nu stă în numărul de funcționalități, ci în disciplina cu care au fost
construite: o singură sursă de adevăr validată strict, o suită de teste de integrare pe infrastructură
reală și o evaluare cantitativă a unei alegeri de arhitectură. Aceste rezultate sunt direct
transferabile: tiparul de instantanee pentru adevăr istoric, modelul de citire derivat și consistent
și metodologia de comparare a soluțiilor de căutare se aplică multor altor domenii (comerț
electronic, sisteme de documente, jurnale de activitate).

\section{Dezvoltări ulterioare}

Arhitectura a fost proiectată pentru extindere, iar schema bazei de date anticipează deja mai multe
direcții (tabelele \code{outbox\_events}, \code{template\_analysis\_results},
\code{coach\_session\_comments} și \code{coach\_template\_assignments} există din prima migrare, încă
neutilizate de cod). Direcțiile imediate, prezentate ca proiectare în Secțiunea~\ref{sec:directii},
sunt: (1) un strat de cache Redis de tip cache-aside pentru marketplace, cu invalidare condusă de
evenimente și protecție împotriva efectului de turmă; (2) un \emph{outbox tranzacțional} pentru
indexare, care rezolvă problema dual-write reutilizând tabela și indexatorul idempotent existente;
(3) o evaluare a calității relevanței (precision@k, nDCG), complementară evaluării de viteză deja
realizate; (4) paginare prin chei (keyset) pentru pagini adânci; (5) limitarea ratei la autentificare
și migrarea la Argon2id; și (6) observabilitate prin trasare distribuită și metrici, partiționarea
temporală a istoricului, o coadă de lucru bazată pe \code{SKIP LOCKED} și vederi materializate pentru
analize. Fiecare dintre acestea extinde, nu remodelează, fundația existentă.
